\documentclass[10pt,a4paper]{amsart}
\usepackage[margin=19mm]{geometry}
\usepackage{amsmath,amssymb,mathtools}
\usepackage[scr=rsfs,cal=euler]{mathalfa}
\usepackage{xfrac}
\usepackage[unicode,hidelinks,bookmarks=false]{hyperref}
\newcommand{\ie}{i.e.\ }
\newcommand{\cf}{cf.\ }
\newcommand{\resp}{respectively\ }
\newcommand{\Subsection}{\subsection}
\providecommand{\nobreakdash}{\nobreak\hbox{-}}
\setlength{\emergencystretch}{2em}
\multlinegap0pt
\usepackage{amssymb}
\usepackage{xcolor}
\usepackage[shortlabels]{enumitem}
\newtheorem{lem}{Lemma}
\newcommand{\R}{\mathbb{R}}
\title{The Multiple Timescales of Gradient Descent on~the~Edge~of~Stability: A~Perturbative~Derivation~of~the~Central~Flow}
\author{Rapha\"el Berthier}
\date{Source: arXiv:2609.01034v1 (CC BY 4.0)}
\begin{document}
\maketitle

\noindent\emph{Source and licence.} The Multiple Timescales of Gradient Descent on~the~Edge~of~Stability: A~Perturbative~Derivation~of~the~Central~Flow. Version arXiv:2609.01034v1 (CC BY 4.0), distributed by arXiv under the Creative Commons Attribution 4.0 International licence, http://creativecommons.org/licenses/by/4.0/, as recorded in the arXiv licence metadata for this exact version and shown on the abstract page https://arxiv.org/abs/2609.01034v1. Full licence notice: \texttt{licenses/arxiv-2609.01034-CC-BY-4.0.txt}.\par\smallskip

\noindent\textit{Editorial scope.} Counted unit: the article's lem lem:secular-general (source0.tex lines 1032--1043). The article's complete original proof of it is delivered with it (source0.tex lines 1045--1060, one \texttt{proof} environment of 16 lines). No mathematics is written, repaired or re-derived: the statement and the proof are the article's own lines, in the article's own notation, and no retained line is substituted or re-flowed.  No further source-local context is needed: every cross-reference of every retained line resolves inside the two delivered spans.  Nothing was omitted from any retained span, so the package carries no omission record.  No retained line contains a citation, so the wrapper carries no bibliography and no bibliography entry.  No retained line contains a figure environment, an image-inclusion command or drawing code, so no image is delivered and the package declares no asset dependency.  Editorial formatting only, in addition: an AMS \texttt{amsart} preamble that restates the article's own declarations and macros used by the retained lines, namely the environments \texttt{lem} on separate counters, together with the standard packages those lines need.  The retained environments print in this document as Lemma 1 in order of appearance, whereas the article prints the same items with the numbering of its own full document.  The complete native source of the article as distributed in the arXiv e-print for this exact version is bundled unchanged as \texttt{sources/209075/source0.tex}.
    \begin{lem}
        \label{lem:secular-general}
        Let $\alpha, \beta \in \R$ and $\gamma = \gamma(t_1) \in \R$ such that $\sum_{t_1=0}^{+\infty} |\gamma(t_1)| < +\infty$. The solutions $\varphi = \varphi(t_1)$ of the difference equation
        \begin{equation*}
            \Delta_{t_1} \varphi + 2 \varphi = \alpha + \beta (-1)^{t_1} + \gamma(t_1)
        \end{equation*}
        are 
        \begin{equation*}
            \varphi = \frac{\alpha}{2} + a (-1)^{t_1} - \beta t_1 (-1)^{t_1} - \sum_{s=0}^{t_1-1} (-1)^{t_1+s} \gamma(s) \, , \qquad a \in \R \, .
        \end{equation*} 
        In particular, there exists a bounded solution if and only if $\beta = 0$.
    \end{lem}

    \begin{proof}
        As $(-1)^{t_1}$ is solution to the homogeneous equation, we seek solutions of the form $\varphi = A(-1)^{t_1}$ for some $A = A(t_1)$. We compute
        \begin{align*}
            \Delta_{t_1} \varphi + 2 \varphi &= \varphi(t_1 + 1) - \varphi(t_1) + 2 \varphi(t_1) = \varphi(t_1 + 1) + \varphi(t_1) \\
            &= A(t_1 + 1)(-1)^{t_1 + 1} + A(t_1)(-1)^{t_1} = -(-1)^{t_1} \Delta_{t_1} A \, .
        \end{align*}
        The function $\varphi = A(-1)^{t_1}$ solves the difference equation if and only if $A$ solves
        \begin{equation*}
            \Delta_{t_1} A = - \alpha (-1)^{t_1} - \beta - (-1)^{t_1} \gamma(t_1) \, .
        \end{equation*}
        The discrete primitives are 
        \begin{equation*}
            A = a + \frac{\alpha}{2} (-1)^{t_1} - \beta t_1 - \sum_{s=0}^{t_1-1} (-1)^s \gamma(s) \, , \qquad a \in \R \, .
        \end{equation*}
        This gives the result.
    \end{proof}

\end{document}
